\guidechapter{Vectors, rank, and a frame you can draw}{The geometry behind the technical vocabulary.}

A vector can be as simple as a pair of numbers, such as $(1,2)$.
Draw it as an arrow from the origin to that point. Adding vectors adds their
coordinates. Multiplying one by a number scales the arrow. A subspace is a
collection closed under these operations: a line through the origin is one
example, and the whole plane is another.

A matrix is a rule for combining input coordinates into output coordinates.
Its columns tell you where the basic input directions go. If every output
lies on a single line, there is only one independent output direction, even
if the matrix has many nonzero entries. That number of independent directions
is its rank.

\explfigure{17-rank-directions}{Count independent directions, not arrows.}{
On the left, $(2,2)$ is twice $(1,1)$: both arrows lie on one line.
On the right, no multiple of $(1,0)$ reaches $(0,1)$: both directions
are needed. These represent ranks one and two respectively.}

The matrices
\[
 \begin{pmatrix}1&2\\1&2\end{pmatrix}
 \quad\hbox{and}\quad
 \begin{pmatrix}1&0\\0&1\end{pmatrix}
\]
have ranks one and two. The first has four nonzero entries but lower rank.
For that first matrix, input $(s,t)$ produces $(s+2t,s+2t)$, always on the
diagonal line. The second can produce every point in the plane. This is why
counting matrix entries would miss the saving sought by the paper.

\paragraph{A frame is a consistent representation rule.}
Take $D$ to rotate every vector a quarter-turn counterclockwise. The logical
value $(1,0)$ is stored as $(0,1)$. The logical value $(0,1)$ is stored as
$(-1,0)$. Adding the stored values gives $(-1,1)$, exactly the rotation of
the logical sum $(1,1)$. Both inputs used the same rule, so addition remained
meaningful.

\explfigure{18-frame-example}{Mixing representations changes the answer.}{
Blue and gold are the two inputs; teal is their sum. Rotating both inputs
rotates the correct sum. Rotating just the first produces $(0,2)$, which
is not the intended rotated result $(-1,1)$. Overlapping arrows in the
right panel deliberately show the two equal directions.}

In the paper, a frame acts on an entire address-indexed array rather than
just a two-number vector. The same law of linearity applies. The later
two-by-two tile diagram gives a concrete array version: rotate the address
grid of each input and then add matching cells. You get the rotated sum.

\paragraph{Why changing frames can be worthwhile.}
One representation may make a particular gate cheap while another makes
the next gate cheap. The algorithm can switch, provided it pays for switching
and maintains the interpretation consistently. Its endpoint frames determine
the intended final operation; its intermediate frames are chosen to keep
the sum of transition costs small. The actual proof uses more sophisticated
operators than rotations, and verifies their residual ranks explicitly.

\takeaway{Rank measures independent directions. A frame is an invertible
representation. Shared frames preserve linear gates; unpriced frame changes
would hide real work.}
